\documentclass[11pt]{article}
\usepackage[a4paper,margin=2.4cm]{geometry}
\usepackage{xeCJK}
\usepackage{amsmath,amssymb}
\usepackage{enumitem}
\usepackage{xcolor}
\usepackage{hyperref}
\setCJKmainfont{Microsoft YaHei}
\setlist{nosep,leftmargin=1.6em}
\pagestyle{plain}

% ── 盒子：手写，只用核心命令 ──────────────────────────────────
% 原本用的是 tcolorbox，但便携 MiKTeX 里缺它的依赖 trimspaces.sty，
% 而自动装包在这台机器上没生效。与其去补依赖，不如不用它——
% 手写盒子的好处是**只依赖 LaTeX 内核**，换任何一台装了 xelatex 的机器都能编。
%
% ⚠️ 这里刻意用 \newcommand 而不是 \newenvironment：
% \newenvironment{...}[2]{begin}{end} 的**结束段不能引用参数**，
% 写了 #1 会报 "Illegal parameter number in definition of \end...".
% 普通命令没有这个限制。
\newcommand{\hwboxout}[3]{%
  \par\medskip\noindent
  \setlength{\fboxsep}{7pt}%
  \fcolorbox{#1}{#2}{\begin{minipage}{0.90\textwidth}#3\end{minipage}}%
  \par\medskip}

\title{\textbf{作业求解结果}}
\author{由 homework-solver 自动生成（求解 $\to$ 自校验 $\to$ 排版）}
\date{\today}
\begin{document}
\maketitle
\vspace{-1.2em}
\hwboxout{gray!70}{gray!8}{\textbf{本文件由流水线自动生成。}每题都经过\textbf{自校验}——
解出的结果会代回原题复算一次；\textbf{自校验未通过的题会被红框标出}，
不会安静地混在正确答案里。}
\vspace{0.2em}

\section*{总览}

\begin{itemize}

  \item 题目总数：\textbf{2}
  \item 成功求解：\textbf{2}
  \item 自校验通过：\textbf{0}
  \item 自校验\textbf{未通过}：\textbf{0}
\end{itemize}

\section*{第 1 题\quad 应用题（一元二次建模）}

\textbf{题目.}\;一个长方形的长比宽多 3 厘米，它的面积是 40 平方厘米。求这个长方形的长和宽。


\textbf{求解过程.}

\begin{enumerate}[label=\arabic*.]

  \item $	ext{设}$ $a$ 为长方形的宽，$b$ 为长方形的长。
  \item $	ext{则有}$ $a + b = 40$ 平方厘米。
\end{enumerate}

\hwboxout{gray!70}{gray!8}{\textbf{答案：}\(通过解这个方程，我们可以找到长方形的宽和长。\)}

\textbf{自校验：}不适用——由模型求解，**无自动校验**；仅供人工复核

\vspace{0.6em}

\section*{第 2 题\quad 古典概型（文字题）}

\textbf{题目.}\;从 1, 2, 3, 4, 5 这五个数字中任取两个不同的数，求这两个数的和为偶数的概率。


\textbf{求解过程.}

\begin{enumerate}[label=\arabic*.]

  \item 第一步：计算从 1, 2, 3, 4, 5 中任取两个不同的数的情况数。
  \item 第二步：计算这些情况数中，两个数的和为偶数的情况数。
  \item 第三步：根据概率公式计算这两个数的和为偶数的概率。
\end{enumerate}

\hwboxout{gray!70}{gray!8}{\textbf{答案：}\(从 1, 2, 3, 4, 5 这五个数字中任取两个不同的数，求这两个数的和为偶数的概率是 $rac{6}{10} = rac{3}{5}$。\)}

\textbf{自校验：}不适用——由模型求解，**无自动校验**；仅供人工复核

\vspace{0.6em}


\end{document}
